% !TeX root = ../main.tex
\subsection{Исследование влияния разрядности цифрового гетеродина}

В данном пункте анализируется зависимость уровня побочных (паразитных) спектральных составляющих на выходе демодулятора от количества бит цифрового гетеродина $n_g$. Моделирование проводится в режиме фиксированной запятой (\texttt{nco\_type = 'fixed'}) для значений $n_g$ вплоть до $n_a + 2 = 15$ бит.

% --- Блок 1: ng = 3 и ng = 6 ---
\begin{figure}[H]
    \centering
    \begin{minipage}[t]{0.48\textwidth}
        \vspace{0pt}
        \centering
        \includegraphics[width=\linewidth]{3_spec_ng3.png}
        \caption{Спектр на выходе DEM ($n_g = 3$)}
        \label{fig:3_ng3}
    \end{minipage}
    \hfill
    \begin{minipage}[t]{0.48\textwidth}
        \vspace{0pt}
        \centering
        \includegraphics[width=\linewidth]{3_spec_ng6.png}
        \caption{Спектр на выходе DEM ($n_g = 6$)}
        \label{fig:3_ng6}
    \end{minipage}
\end{figure}

Как видно из рисунков \ref{fig:3_ng3} и \ref{fig:3_ng6}, при крайне малой разрядности гетеродина ($n_g = 3$ и $n_g = 6$ бит) ошибка квантования формируемых синусоидальных и косинусоидальных отсчетов весьма велика. Это приводит к нелинейным искажениям и появлению мощных побочных гармоник в спектре выходного сигнала. Данные паразитные составляющие существенно сокращают динамический диапазон устройства, свободный от искажений, и могут привести к ложному обнаружению сигналов.

% --- Блок 2: ng = 10 и ng = 13 ---
\begin{figure}[H]
    \centering
    \begin{minipage}[t]{0.48\textwidth}
        \vspace{0pt}
        \centering
        \includegraphics[width=\linewidth]{3_spec_ng10.png}
        \caption{Спектр на выходе DEM ($n_g = 10$)}
        \label{fig:3_ng10}
    \end{minipage}
    \hfill
    \begin{minipage}[t]{0.48\textwidth}
        \vspace{0pt}
        \centering
        \includegraphics[width=\linewidth]{3_spec_ng13.png}
        \caption{Спектр на выходе DEM ($n_g = 13$)}
        \label{fig:3_ng13}
    \end{minipage}
\end{figure}

С увеличением разрядности гетеродина до $n_g = 10$ и $n_g = 13$ бит (рисунки \ref{fig:3_ng10} и \ref{fig:3_ng13}) уровень шума квантования монотонно снижается. Паразитные спектральные пики становятся значительно слабее, так как форма синтезируемого опорного сигнала приближается к идеальной гармонике. 

% --- Блок 3: ng = 15 ---
\begin{figure}[H]
    \centering
    \begin{minipage}[t]{0.48\textwidth}
        \vspace{0pt}
        \centering
        \includegraphics[width=\linewidth]{3_spec_ng15.png}
        \caption{Спектр на выходе DEM ($n_g = 15$)}
        \label{fig:3_ng15}
    \end{minipage}
\end{figure}

При достижении значения $n_g = 15$, что соответствует рекомендованному правилу $n_g = n_a + 2$ (рисунок \ref{fig:3_ng15}), мощность шумов квантования гетеродина становится пренебрежимо малой по сравнению с шумами квантования самого АЦП. Дальнейшее увеличение разрядности ЦГ нецелесообразно, так как общий динамический диапазон тракта будет ограничен характеристиками аналого-цифрового преобразователя.
